\appendix
\section*{ПРИЛОЖЕНИЯ}
\addcontentsline{toc}{section}{ПРИЛОЖЕНИЯ}

% По ГОСТ 7.32-2017 (п. 6.14.3) таблицы и рисунки в приложении
% нумеруются буквой приложения и порядковым номером: А.1, А.2, ...
\renewcommand{\thetable}{А.\arabic{table}}
\renewcommand{\thefigure}{А.\arabic{figure}}
\setcounter{table}{0}
\setcounter{figure}{0}

% =====================================================================
\section*{Приложение А. Полная таблица результатов эксперимента E1}
\addcontentsline{toc}{subsection}{Приложение А. Полная таблица результатов E1}
\label{app:E1-raw}

В \tabref{tab:appE1} приведена полная сводка результатов 22~запусков
эксперимента E1. Для каждой комбинации (модель, seed) зафиксированы:
оптимальный порог $\tau$ (под target FPR\,=\,0{,}01 на val),
температура калибровки, ECE до/после калибровки, Precision, Recall,
F1, FPR, MCC, ROC-AUC, PR-AUC с bootstrap-CI~95\,\%, latency на одно
окно.

{
\renewcommand{\arraystretch}{1.15}
\setlength{\tabcolsep}{4pt}
\footnotesize
\begin{longtable}{@{}lrrrrrrrr@{}}
\caption{Полная сводка результатов эксперимента E1 (22 запуска)}
\label{tab:appE1}\\
\toprule
\textbf{Модель} & \textbf{seed} & \textbf{$\tau$} & \textbf{F1} &
\textbf{PR-AUC} & \textbf{Prec.} & \textbf{Rec.} & \textbf{FPR} &
\textbf{Lat., мс} \\
\midrule
\endfirsthead
\multicolumn{9}{@{}l@{}}{\textit{Продолжение таблицы~\ref{tab:appE1}}}\\
\toprule
\textbf{Модель} & \textbf{seed} & \textbf{$\tau$} & \textbf{F1} &
\textbf{PR-AUC} & \textbf{Prec.} & \textbf{Rec.} & \textbf{FPR} &
\textbf{Lat., мс} \\
\midrule
\endhead
\bottomrule
\endfoot
logistic\_regression & 42   & 0.502 & 0.847 & 0.986 & 0.991 & 0.739 & 0.014 & 0.11  \\
logistic\_regression & 123  & 0.502 & 0.847 & 0.986 & 0.991 & 0.739 & 0.014 & 0.12  \\
logistic\_regression & 2024 & 0.502 & 0.847 & 0.986 & 0.991 & 0.739 & 0.014 & 0.08  \\
random\_forest       & 42   & 0.512 & 0.917 & 0.993 & 0.992 & 0.853 & 0.014 & 54.83 \\
random\_forest       & 123  & 0.515 & 0.917 & 0.993 & 0.991 & 0.852 & 0.016 & 52.74 \\
random\_forest       & 2024 & 0.510 & 0.917 & 0.993 & 0.993 & 0.852 & 0.013 & 46.95 \\
xgboost              & 42   & 0.498 & 0.919 & 0.993 & 0.990 & 0.857 & 0.019 & 0.46  \\
xgboost              & 123  & 0.498 & 0.919 & 0.993 & 0.990 & 0.857 & 0.019 & 0.74  \\
xgboost              & 2024 & 0.498 & 0.919 & 0.993 & 0.990 & 0.857 & 0.019 & 0.87  \\
isolation\_forest    & 42   & 0.952 & 0.121 & 0.872 & 0.988 & 0.064 & 0.001 & 21.02 \\
isolation\_forest    & 123  & 0.910 & 0.453 & 0.879 & 0.985 & 0.296 & 0.009 & 21.34 \\
isolation\_forest    & 2024 & 0.940 & 0.198 & 0.872 & 0.987 & 0.110 & 0.004 & 21.57 \\
cnn\_lstm            & 42   & 0.675 & 0.845 & 0.910 & 0.992 & 0.561 & 0.033 & 0.95  \\
cnn\_lstm            & 123  & 0.682 & 0.722 & 0.975 & 0.987 & 0.569 & 0.013 & 1.12  \\
cnn\_lstm            & 2024 & 0.680 & 0.743 & 0.980 & 0.987 & 0.595 & 0.012 & 0.82  \\
mlp                  & 42   & 0.500 & 0.809 & 0.715 & 0.679 & 1.000 & 1.000 & 0.18  \\
cnn1d                & 42   & 0.700 & 0.674 & 0.979 & 0.993 & 0.510 & 0.008 & 0.69  \\
lstm                 & 42   & 0.690 & 0.821 & 0.985 & 0.991 & 0.701 & 0.013 & 0.82  \\
gru                  & 42   & 0.685 & 0.821 & 0.985 & 0.988 & 0.703 & 0.018 & 3.39  \\
bilstm               & 42   & 0.695 & 0.842 & 0.988 & 0.992 & 0.731 & 0.012 & 1.84  \\
tcn                  & 42   & 0.692 & 0.844 & 0.986 & 0.992 & 0.735 & 0.013 & 4.20  \\
transformer          & 42   & 0.700 & 0.773 & 0.981 & 0.989 & 0.635 & 0.016 & 1.44  \\
\end{longtable}
}

\noindent
Сырая таблица в формате CSV (\verb|E1_results.csv|) поставляется в
составе артефактов программного прототипа.

% =====================================================================
\section*{Приложение Б. Структура программного прототипа}
\addcontentsline{toc}{subsection}{Приложение Б. Структура программного прототипа}
\label{app:project}

Прототип реализован как Python-пакет \texttt{diploma\_nids} с
подпакетами \texttt{data, models, training, eval, attacker, inference,
ui, utils}. Полное описание модулей и тестов приведено на
этапе~\ref{ch:impl} настоящего отчёта.

\textbf{Запуск экспериментов} (из корневой директории программного
прототипа):

\begin{verbatim}
make install-dev
python scripts/02_build_dataset.py \
    --config configs/data/unsw_nb15.yaml \
    --preprocess configs/preprocess/default.yaml \
    --strategy official
python scripts/10_run_experiment_E1.py \
    --train configs/train/full.yaml --seeds 42 123 2024
python scripts/11_run_experiment_E2_error_analysis.py
python scripts/12_run_experiment_E4_attacker.py
python scripts/13_run_experiment_E5_drift.py
python scripts/20_build_thesis_tables.py
python scripts/22_parse_training_logs.py
python scripts/21_plot_training_history.py
python scripts/30_demo_run.py
\end{verbatim}

% =====================================================================
\section*{Приложение В. Конфигурации экспериментов}
\addcontentsline{toc}{subsection}{Приложение В. Конфигурации экспериментов}
\label{app:configs}

Все стадии pipeline параметризуются YAML-файлами в каталоге
\verb|configs/| программного прототипа. Ключевые группы конфигов:

\begin{itemize}
    \item \verb|configs/data/| --- схемы UNSW-NB15 и CICIDS2017
    (mapping общих признаков для cross-evaluation);
    \item \verb|configs/preprocess/default.yaml| --- параметры
    очистки, encoding, log/robust scaling, windowing, splits,
    балансировки;
    \item \verb|configs/models/| --- архитектуры 10 DL-моделей
    (\verb|cnn_lstm.yaml| --- proposed) и параметры 5 классических
    моделей (\verb|classical.yaml|);
    \item \verb|configs/train/| --- параметры обучения (\verb|full.yaml|
    для полного режима, \verb|smoke.yaml| для smoke-тестов);
    \item \verb|configs/eval/default.yaml| --- метрики, калибровка,
    bootstrap CI, drift-monitoring;
    \item \verb|configs/attacker/| --- политика FSM (\verb|policy.yaml|),
    7 шаблонов атак, drift-инжектор (\verb|drift.yaml|);
    \item \verb|configs/pipeline/realtime.yaml| --- конфигурация
    realtime-стенда FastAPI + Streamlit + agent.
\end{itemize}

Каждый запуск экспериментов фиксирует: версию кода (commit hash),
набор seed, гиперпараметры, версии зависимостей, полное содержимое
использованных конфигов в виде attached-артефактов.

% Метка последнего рисунка/таблицы для подсчёта в реферате
\label{lastfig}
\label{lasttab}
